\documentclass[11pt]{article}

\usepackage[margin=1in]{geometry}
\usepackage{amsmath}
\usepackage{amssymb}
\usepackage{booktabs}
\usepackage{array}
\usepackage{enumitem}
\usepackage[colorlinks=true,linkcolor=blue,urlcolor=blue]{hyperref}

\setlength{\parindent}{0pt}
\setlength{\parskip}{0.55em}

\newcommand{\pts}[1]{\hspace{0.35em}\textnormal{\textbf{[#1 points]}}}
\newenvironment{choices}{\begin{enumerate}[label=\Alph*.,leftmargin=2.2em,itemsep=0.12em,topsep=0.12em]}{\end{enumerate}}

\title{S181M116 Derivatives and Alternative Investments\\
Final Exam, Variant B: Multiple Choice}
\author{Time allowed: 1 hour 30 minutes \qquad Total: 80 points}
\date{}

\begin{document}
\maketitle

\textbf{Instructions.} Answer all questions. Each item has exactly one best
answer. Unless stated otherwise, ignore transaction costs, taxes, bid-ask
spreads, and default risk. Use continuous compounding. Round monetary answers to
two decimals, probabilities to four decimals, and hedge ratios to four decimals.

\textbf{Scoring.} There are 20 multiple-choice items. Each item is worth 4
points, for a total of 80 points.

\textbf{Coverage.} This exam is based on the material covered in class:
derivative basics, market structure, forwards and futures, hedging and basis
risk, forward/futures pricing, options, put-call parity, binomial trees, and
option no-arbitrage restrictions. Lecture 3 and Lecture 5 material are excluded.

\section*{Question 1: Concepts, Payoffs, and Market Structure \pts{16}}

This question is on the easier side.

\begin{enumerate}
    \item A contract pays $\max(90-S_T,0)$ in six months. Which statement is
    correct? \pts{4}
    \begin{choices}
        \item It is not a derivative because it has no initial premium.
        \item It is a derivative whose underlying variable is the asset price $S_T$.
        \item It is a bond because the payoff is bounded above.
        \item It is a futures contract because the payoff is linear.
    \end{choices}

    \item A European call and a European put both have strike $K=75$. The call
    premium is 6 and the put premium is 4. If $S_T=88$, what are the buyer's
    call payoff, call profit, put payoff, and put profit? \pts{4}
    \begin{choices}
        \item $13,\;7,\;0,\;-4$
        \item $13,\;13,\;0,\;0$
        \item $0,\;-6,\;13,\;9$
        \item $7,\;13,\;-4,\;0$
    \end{choices}

    \item Which statement best describes an OTC derivative relative to an
    exchange-traded derivative? \pts{4}
    \begin{choices}
        \item OTC contracts are more standardized and always centrally cleared.
        \item OTC contracts are customizable, but can create more counterparty,
        valuation, and liquidity risk.
        \item OTC contracts have no counterparty risk because they are private.
        \item OTC contracts are always cheaper because they have no margin.
    \end{choices}

    \item Why can margin create liquidity pressure? \pts{4}
    \begin{choices}
        \item Because daily settlement can require losing positions to post cash
        before the underlying exposure has generated offsetting cash flow.
        \item Because margin is the purchase price of the futures contract.
        \item Because variation margin removes all price risk.
        \item Because margin calls occur only when the hedge is profitable.
    \end{choices}
\end{enumerate}

\section*{Question 2: Futures Settlement and Forward Pricing \pts{16}}

This question is on the easier side.

\begin{enumerate}
    \item A trader is \textbf{short} eight copper futures contracts. Each
    contract is for 25,000 pounds. Yesterday's settlement price was USD 4.12 per
    pound and today's settlement price is USD 4.05 per pound. What is today's
    futures profit or loss? \pts{4}
    \begin{choices}
        \item USD 14,000 gain
        \item USD 14,000 loss
        \item USD 1,400 gain
        \item USD 1,400 loss
    \end{choices}

    \item The trader's margin account balance before settlement was USD 38,000.
    The maintenance margin is USD 30,000 and the initial margin is USD 38,000.
    What is the margin account balance after settlement, and is there a margin
    call? \pts{4}
    \begin{choices}
        \item USD 24,000; yes, deposit USD 14,000.
        \item USD 30,000; no margin call.
        \item USD 52,000; no margin call.
        \item USD 52,000; yes, deposit USD 8,000.
    \end{choices}

    \item A stock has current price $S_0=120$. The present value of known cash
    dividends during a six-month forward contract is $I=2.40$. The continuously
    compounded risk-free rate is 3.5\% per year. What is the no-arbitrage
    forward price? \pts{4}
    \begin{choices}
        \item 117.60
        \item 119.68
        \item 122.10
        \item 124.25
    \end{choices}

    \item A dealer quotes a forward price of 121 for the same contract. Which
    statement is correct? \pts{4}
    \begin{choices}
        \item The quote is too low; short the asset and long the forward.
        \item The quote is too high; buy the asset, short the forward, and lock
        in about 1.32 per share.
        \item The quote is exactly no-arbitrage.
        \item The quote is too high; short the asset, invest proceeds, and long
        the forward.
    \end{choices}
\end{enumerate}

\section*{Question 3: Cross Hedging, Basis Risk, and Effective Prices \pts{16}}

This question is on the harder side.

A manufacturer will buy 1,250,000 pounds of copper in three months. It cross
hedges with a related metals futures contract. Each futures contract covers
25,000 pounds. Historical data give
\[
    \rho=0.72,\qquad \sigma_S=0.18,\qquad \sigma_F=0.16.
\]
The hedge is opened at $F_1=3.80$. At hedge close, the manufacturer buys copper
at $S_2=4.12$ and closes futures at $F_2=4.05$.

\begin{enumerate}
    \item What hedge direction should the manufacturer use? \pts{4}
    \begin{choices}
        \item Long futures, because it will buy the commodity later and wants
        protection against rising prices.
        \item Short futures, because it will buy the commodity later.
        \item Long futures only if the expected basis is negative.
        \item No hedge is possible because the contract is a cross hedge.
    \end{choices}

    \item What is the minimum-variance hedge ratio? \pts{4}
    \begin{choices}
        \item 0.6400
        \item 0.7200
        \item 0.8100
        \item 1.1250
    \end{choices}

    \item What number of futures contracts is implied by the hedge ratio before
    rounding? \pts{4}
    \begin{choices}
        \item 36.0
        \item 40.5
        \item 50.0
        \item 56.25
    \end{choices}

    \item What is the effective purchase price per pound, and how does the
    realized basis compare with an expected basis of 0.03? \pts{4}
    \begin{choices}
        \item 3.87; the realized basis is worse for the buyer.
        \item 3.87; the realized basis is better for the buyer.
        \item 4.37; the realized basis is worse for the buyer.
        \item 4.37; the realized basis is better for the buyer.
    \end{choices}
\end{enumerate}

\section*{Question 4: Option No-Arbitrage and Binomial Pricing \pts{16}}

This question is on the harder side.

Consider European options on a non-dividend-paying stock. The current stock
price is $S_0=48$, the strike is $K=50$, the continuously compounded risk-free
rate is $r=5\%$, and maturity is $T=0.25$ years. For the parity questions, a
European call trades at $c=2.60$. Separately, for the binomial model questions,
use $u=1.18$ and $d=0.88$.

\begin{enumerate}
    \item What is the no-arbitrage European put price implied by put-call
    parity? \pts{4}
    \begin{choices}
        \item 1.22
        \item 2.60
        \item 3.20
        \item 3.98
    \end{choices}

    \item Suppose the market put price is $p=3.20$. Which arbitrage statement is
    correct? \pts{4}
    \begin{choices}
        \item The protective put is expensive; sell it and buy the fiduciary call.
        \item The fiduciary call is expensive; sell it and buy the protective put.
        \item Both packages are fairly priced.
        \item There is no arbitrage because call and put prices are both positive.
    \end{choices}

    \item In the one-step tree, what are the up-state and down-state stock
    prices and call payoffs? \pts{4}
    \begin{choices}
        \item $S_u=56.64$, $S_d=42.24$; payoffs $6.64$ and $0$.
        \item $S_u=56.64$, $S_d=42.24$; payoffs $0$ and $7.76$.
        \item $S_u=59.00$, $S_d=44.00$; payoffs $9.00$ and $0$.
        \item $S_u=48.18$, $S_d=47.88$; payoffs $0$ and $0$.
    \end{choices}

    \item Which pair is closest to the risk-neutral probability and binomial call
    value? \pts{4}
    \begin{choices}
        \item $q=0.4419$, call value $2.90$
        \item $q=0.4419$, call value $6.64$
        \item $q=0.5581$, call value $2.60$
        \item $q=0.5581$, call value $3.67$
    \end{choices}
\end{enumerate}

\section*{Question 5: Option Bounds, Early Exercise, and Strategy Logic \pts{16}}

This question is on the harder side.

Consider options on a non-dividend-paying stock with current price $S_0=70$,
strike $K=68$, risk-free rate $r=4.5\%$ with continuous compounding, and maturity
$T=0.75$ years.

\begin{enumerate}
    \item What is the lower bound for the European call, and what does it imply
    if the call is quoted at $c=3.80$? \pts{4}
    \begin{choices}
        \item Lower bound about 0; no arbitrage.
        \item Lower bound about 1.74; no arbitrage.
        \item Lower bound about 4.26; the call violates the lower bound.
        \item Lower bound about 6.00; the call is fairly priced.
    \end{choices}

    \item What is the lower bound for the European put, and what does it imply
    if the put is quoted at $p=0.60$? \pts{4}
    \begin{choices}
        \item Lower bound 0; no lower-bound arbitrage from this quote.
        \item Lower bound 1.74; the put violates the lower bound.
        \item Lower bound 4.26; the put violates the lower bound.
        \item Lower bound 68; the put violates the lower bound.
    \end{choices}

    \item If a European call with this strike and maturity trades at $c=6.00$,
    what is the no-arbitrage European put price from put-call parity? \pts{4}
    \begin{choices}
        \item 0.00
        \item 1.74
        \item 3.80
        \item 6.00
    \end{choices}

    \item Which statement is correct? \pts{4}
    \begin{choices}
        \item An American put price of $P=1.50$ is consistent if the European put
        value from parity is about 1.74.
        \item A covered call is risk-free because the writer owns the stock.
        \item An American put must be worth at least as much as the otherwise
        identical European put, and a covered call still has downside stock risk.
        \item A covered call has unlimited upside because the writer owns the stock.
    \end{choices}
\end{enumerate}

\end{document}
